% Responsable: Danela · Cierre (domingo 15, hasta 21.00)
% Declaración de uso de IA de los tres integrantes.

% Chavarria Tahiny 
\section*{Declaraci\'on de uso de inteligencia artificial: Tahiny Mel Chavarr\'ia Cuenca}

Se us\'o Claude, asistente de inteligencia artificial de Anthropic, el 12 de noviembre de 2026, solo en mis partes del trabajo. La IA se us\'o para:

\begin{itemize}
  \item redactar un borrador de las secciones 2.3 (requisitos de sistema y de software) y 2.4 (modelo de calidad ISO/IEC 25010:2023);
  \item proponer el diseño de la tabla de 15 requisitos de El Balc\'on Restaurant, con su tipo, redacci\'on EARS, prioridad MoSCoW y fuente;
  \item redactar un borrador de mis dos conclusiones;
  \item preparar las entradas BibTeX de las referencias y la estructura de los archivos LaTeX.
\end{itemize}

% Arteaga Danela
\section*{Declaración de uso de inteligencia artificial: Danela Dayana Arteaga Alava}

Se usó Claude, asistente de inteligencia artificial de Anthropic, entre el 12 y el 15 de noviembre de 2026, en mis partes del trabajo, la IA se usó para:

\begin{itemize}
  \item Proponer borradores de la Sección 0, de las secciones 2.1 y 2.2, de la evaluación de los requisitos y de la aplicación del proceso, y una primera versión del mapa conceptual.
  \item Verificar los datos de las referencias contra fuentes oficiales y comprobar las cláusulas citadas.
  \item Organizar el reparto de tareas, la estructura de archivos LaTeX del grupo y el README.
  \item Corregir faltas ortográficas y ambigüedad.
  \item Revisar si hacía falta algo en el repositorio, básicamente fue de ayuda en revisar que todo en el repositorio esté correctamente subido.
\end{itemize}

% Naranjo Anthony
\section*{Declaración de uso de inteligencia artificial: Naranjo Alfonzo Anthony Alexander}

Se usó ChatGPT, asistente de inteligencia artificial de OpenAI, entre el 12 y el 14 de noviembre de 2026, únicamente como apoyo en mis partes asignadas de la práctica. La IA se usó para:

\begin{itemize}
  \item orientar la redacción de la sección 2.5 sobre tipos de sistemas y su aplicación al Sistema de gestión para El Balcón Restaurant con IA;
  \item apoyar la elaboración de la sección 2.6 sobre el proceso de requisitos, modelos de proceso y framework de Pohl;
  \item estructurar el límite del sistema, la declaración de visión y las ocho perspectivas de contexto;
  \item organizar la identificación de los actores del proceso de requisitos y el mapeo de las actividades centrales y transversales del framework de Pohl;
  \item apoyar la redacción y revisión de mis dos conclusiones;
  \item revisar la estructura LaTeX, las citas y la coherencia de mis archivos antes de realizar los commits en GitHub.
\end{itemize}

El contenido generado con apoyo de IA fue revisado, adaptado y comprendido antes de ser incorporado al informe.
